% ============================================================
% DOCUMENTCLASS
% ============================================================
\documentclass[12pt,oneside]{book}

% ============================================================
% METADATOS
% ============================================================
\newcommand{\universidad}{Universidad de Buenos Aires (UBA)}
\newcommand{\facultad}{Facultad de Derecho}
\newcommand{\materia}{Derecho de la Integración}
\newcommand{\titulo}{Apuntes de Derecho de la Integración}
\newcommand{\subtitulo}{Mercosur · Unión Europea · Derechos humanos en la integración}
\newcommand{\autor}{Isaac Edgar Camacho Ocampo}
\newcommand{\anio}{2026}

% ============================================================
% PAQUETES
% ============================================================
\usepackage[utf8]{inputenc}
\usepackage[T1]{fontenc}
\usepackage[spanish,es-nodecimaldot]{babel}

\usepackage[
    top=2.5cm,
    bottom=2.5cm,
    left=3cm,
    right=2.5cm,
    headheight=15pt
]{geometry}

\usepackage{amsmath}
\usepackage{amssymb}
\usepackage{mathtools}

\usepackage{graphicx}
\usepackage{float}
\usepackage{caption}
\usepackage{subcaption}

\usepackage{booktabs}
\usepackage{array}
\usepackage{longtable}
\usepackage{tabularx}

\usepackage{enumitem}

\usepackage{xcolor}
\usepackage[most]{tcolorbox}

\usepackage{fancyhdr}
\usepackage{ragged2e}

\usepackage[
    colorlinks=true,
    linkcolor=blue!50!black,
    citecolor=green!40!black,
    urlcolor=blue!60!black,
    pdftitle={\titulo},
    pdfauthor={\autor},
    pdfsubject={Apuntes universitarios},
    bookmarks=true
]{hyperref}

% ============================================================
% CONFIGURACIÓN
% ============================================================
\graphicspath{
    {./img/}
    {./graficos/}
    {./figuras/}
}

\setcounter{tocdepth}{2}

\setlength{\parindent}{0pt}
\setlength{\parskip}{0.5em}

\setlist[itemize]{
    topsep=0.3em,
    itemsep=0.2em
}

\setlist[enumerate]{
    topsep=0.3em,
    itemsep=0.2em
}

\clubpenalty=10000
\widowpenalty=10000

\pagestyle{fancy}
\fancyhf{}
\fancyhead[L]{\nouppercase{\leftmark}}
\fancyhead[R]{\materia}
\fancyfoot[C]{\thepage}
\renewcommand{\headrulewidth}{0.4pt}
\renewcommand{\footrulewidth}{0pt}

\fancypagestyle{plain}{
    \fancyhf{}
    \fancyfoot[C]{\thepage}
    \renewcommand{\headrulewidth}{0pt}
}

% ============================================================
% COMANDOS
% ============================================================
\newcommand{\abs}[1]{\left|#1\right|}
\newcommand{\norm}[1]{\left\lVert#1\right\rVert}
\newcommand{\vect}[1]{\mathbf{#1}}
\newcommand{\deriv}[2]{\frac{d#1}{d#2}}
\newcommand{\pderiv}[2]{\frac{\partial#1}{\partial#2}}

% ============================================================
% ENTORNOS / CAJAS
% ============================================================
\newtcolorbox{definition}[1]{
    enhanced,
    breakable,
    title={Definición: #1},
    fonttitle=\bfseries,
    colback=blue!3,
    colframe=blue!50!black,
    boxrule=0.8pt,
    arc=2mm
}

\newtcolorbox{important}{
    enhanced,
    breakable,
    title={Importante},
    fonttitle=\bfseries,
    colback=yellow!8,
    colframe=orange!70!black,
    boxrule=0.8pt,
    arc=2mm
}

\newtcolorbox{example}[1]{
    enhanced,
    breakable,
    title={Ejemplo: #1},
    fonttitle=\bfseries,
    colback=green!4,
    colframe=green!40!black,
    boxrule=0.8pt,
    arc=2mm
}

\newtcolorbox{formula}[1]{
    enhanced,
    breakable,
    title={Fórmula / Regla: #1},
    fonttitle=\bfseries,
    colback=purple!4,
    colframe=purple!50!black,
    boxrule=0.8pt,
    arc=2mm
}

\newtcolorbox{exercise}[1]{
    enhanced,
    breakable,
    title={Ejercicio: #1},
    fonttitle=\bfseries,
    colback=gray!5,
    colframe=gray!60!black,
    boxrule=0.8pt,
    arc=2mm
}

\newtcolorbox{note}{
    enhanced,
    breakable,
    title={Nota},
    fonttitle=\bfseries,
    colback=white,
    colframe=gray!50,
    boxrule=0.6pt,
    arc=2mm
}

% ============================================================
% CONFIGURACIÓN FINAL
% ============================================================
\renewcommand{\arraystretch}{1.25}

% ============================================================
% DOCUMENT
% ============================================================
\begin{document}

% ------------------------------------------------------------
% PORTADA
% ------------------------------------------------------------
\begin{titlepage}

\thispagestyle{empty}

\begin{center}

\vspace*{1cm}

{\large \textsc{\universidad}}

\vspace{0.2cm}

{\large \textsc{\facultad}}

\vspace{3cm}

{\Huge \textbf{\materia}}

\vspace{1cm}

{\LARGE \subtitulo}

\vspace{0.8cm}

\textit{Comprender los conceptos es el primer paso para convertir el estudio en conocimiento.}

\vspace{1.2cm}

\rule{0.70\textwidth}{0.5pt}

\vspace{0.5cm}

{\large Apuntes de estudio}

\vspace{0.5cm}

\rule{0.70\textwidth}{0.5pt}

\vfill

{\large \autor}

\vspace{0.3cm}

{\large \anio}

\end{center}

\vfill

\begin{flushleft}
\textit{Este es un modesto aporte a los estudiantes de ``Derecho de la Integración'' no pretende sustituir el material oficial de la materia.}
\end{flushleft}

\end{titlepage}

% ------------------------------------------------------------
% ÍNDICE
% ------------------------------------------------------------
\tableofcontents

% ============================================================
% CONTENIDO
% ============================================================

% ------------------------------------------------------------
\chapter[Aclaraciones: UE y Mercosur]{Aclaraciones conceptuales: Unión Europea y Mercosur}
% ------------------------------------------------------------

\section{Presupuesto en la Unión Europea}

Para determinar quién elabora, fija, aprueba y controla el presupuesto en la Unión Europea debe identificarse qué órgano cumple la función ejecutiva y qué órgano cumple la función de aprobación y control. No corresponde responder con la figura del ministro de Economía, que en este esquema no tiene voz ni voto.

El \textbf{Poder Ejecutivo} es el que representa a la Unión Europea, ejecuta y elabora. La \textbf{Comisión}, a través de los comisarios, es la que lleva adelante y ejecuta todo el plan; elabora el presupuesto, que luego aprueba el Parlamento. El \textbf{Parlamento} es quien aprueba el presupuesto, lo controla y sanciona a los comisarios.

% TODO: contenido incompleto o ambiguo en las notas originales (la exposición menciona primero al Consejo de Ministros como el que ejecuta y elabora, comparándolo con el Senado en cuanto representa a los Estados, y luego precisa que la Comisión es la que lleva adelante y ejecuta el plan)
\begin{note}
En la exposición se nombra primero al Consejo de Ministros como órgano que ejecuta y elabora, comparable al Senado por representar a los Estados, y a continuación se precisa que quien lleva adelante y ejecuta todo el plan es la Comisión. Aquí se consigna la atribución final: el Ejecutivo (la Comisión) elabora y ejecuta; el Parlamento aprueba y controla.
\end{note}

\begin{important}
\begin{itemize}
    \item \textbf{Elabora y ejecuta el plan:} el Ejecutivo (la Comisión, a través de los comisarios).
    \item \textbf{Aprueba y controla el presupuesto:} el Parlamento.
    \item \textbf{Sanciona a los comisarios:} el Parlamento.
\end{itemize}
\end{important}

Las preguntas formuladas con una acción (por ejemplo, ``¿quién legisla tal cosa?'' o ``¿quién juzga tal cosa?'') conducen a identificar el órgano que cumple esa función.

\section{Representación en el Parlamento y en el Consejo}

El \textbf{Parlamento} es el órgano que representa a la población. Por ese motivo, en el Parlamento hay un voto por cada representante y varía la composición y la representación según la proporcionalidad. Por la misma razón, no corresponde exigir al Consejo un porcentaje de población representada (65\,\%).
% TODO: contenido incompleto o ambiguo en las notas originales (referencia al 65 % de la población)

El \textbf{Consejo} se constituye con un representante por cada Estado y un voto por cada Estado.

Esta diferencia en la representación debe tenerse presente al responder, y exige cuidar los términos empleados.

\section{Acceso al tribunal: alternativo y obligatorio}

El carácter \textbf{alternativo} del Tribunal Permanente de Revisión del Mercosur se refiere al \textbf{acceso} al tribunal, que no es obligatorio. No significa que la sentencia no sea obligatoria: es posible acudir al Tribunal Permanente o elegir otro mecanismo de solución de controversias, por ejemplo, la Organización Mundial de Comercio.

En la Unión Europea, en cambio, el único que puede interpretar la aplicación del derecho comunitario es el tribunal. En todo lo relativo a la validez, interpretación y aplicación del derecho comunitario el acceso es obligatorio: hay que acudir a consultarle y acatar lo que diga. En el Mercosur el sistema de solución de conflictos no es obligatorio, sino alternativo: se puede recurrir a este sistema o a otro.

\begin{important}
\begin{itemize}
    \item \textbf{Mercosur (Tribunal Permanente de Revisión):} el acceso es alternativo; puede elegirse otro foro (por ejemplo, la Organización Mundial de Comercio). La sentencia sí es obligatoria.
    \item \textbf{Unión Europea (Tribunal):} el acceso es obligatorio para la validez, interpretación y aplicación del derecho comunitario, y hay que acatar lo que el tribunal diga.
\end{itemize}
\end{important}

\subsection{Cuadro comparativo}

\begin{table}[H]
\centering
\small
\begin{tabularx}{\textwidth}{
    >{\raggedright\arraybackslash}p{0.19\textwidth}
    >{\raggedright\arraybackslash}X
    >{\raggedright\arraybackslash}X
}
\toprule
\textbf{Aspecto} &
\textbf{Mercosur (Tribunal Permanente de Revisión)} &
\textbf{Unión Europea (Tribunal)} \\
\midrule

\textbf{Naturaleza del sistema} &
Sistema mixto: mecanismos de autocomposición y de heterocomposición; incluye un arbitraje ad hoc, ``ni siquiera un arbitraje permanente''. &
Sistema judicial: una instancia judicial en la que se resuelve. \\

\textbf{Acceso} &
Alternativo: se puede ir a este mecanismo o elegir otro (por ejemplo, la Organización Mundial de Comercio). &
Obligatorio: el tribunal es el único que puede interpretar y aplicar el derecho comunitario (validez, interpretación y aplicación). \\

\textbf{Mecanismo de consulta} &
Opinión consultiva: no es necesario hacerle caso. &
Cuestión prejudicial: hay que hacerle caso; es obligatoria. \\

\textbf{Alcance de lo ``alternativo''} &
Lo alternativo es el acceso, no la sentencia: la sentencia sí es obligatoria. &
No hay alternativa: hay que acudir y acatar lo que el tribunal diga. \\

\bottomrule
\end{tabularx}
\end{table}

\section{Derecho originario: Tratado de Asunción y Protocolo de Olivos}

El \textbf{Tratado de Asunción} es derecho originario porque crea el Mercosur: es el tratado que le dio origen.

El \textbf{Protocolo de Olivos} también podría considerarse derecho originario, pero no creó el Mercosur como sí lo hizo el tratado que le dio origen; su consideración como derecho originario puede depender de las expresiones empleadas y de cómo se redacte la respuesta.

\begin{important}
No debe confundirse el \textbf{Protocolo de Olivos} con el ``Pacto de Olivos'': este último es ``la reunión'' y no tiene relación con el tema.
\end{important}

% ------------------------------------------------------------
\chapter[Criterios de estudio y comparación]{Criterios para el estudio y la comparación de sistemas}
% ------------------------------------------------------------

\section{Contenidos de la parte general}

Para la parte general se indicó elaborar, por cada \textbf{tratado fundacional}, un cuadro que consigne:

\begin{itemize}
    \item qué \textbf{área de integración} es y qué características tiene;
    \item si es \textbf{intergubernamental} o \textbf{supranacional};
    \item si es una \textbf{unión aduanera} o una \textbf{zona de libre comercio};
    \item si las decisiones se resuelven por \textbf{consenso} o por \textbf{órganos supranacionales}.
\end{itemize}

Más allá de esas características generales (tipo de integración, resolución por consensos o por órganos supranacionales), no se exigen otras funciones.

\begin{example}{Comunidad Andina}
La Comunidad Andina dispone de órganos supranacionales: sigue el ritmo de la Unión Europea, tiene órganos supranacionales y un derecho que ella misma declara autónomo; es decir, se le dan esas mismas características.
\end{example}

Una pregunta puede requerir conocer qué países pertenecen a un proceso de integración (por ejemplo, ubicar a México dentro de Centroamérica); la respuesta puede alcanzarse por lógica y cultura general.

\section{Lectura de consignas y precisión de los términos}

\begin{itemize}
    \item Debe leerse con atención la consigna: puede haber \textbf{más de una opción correcta}.
    \item Hay que cuidar los términos: \textbf{alternativo} se refiere al acceso al tribunal y no a la obligatoriedad de la sentencia.
    \item Una respuesta pedida de forma breve (cinco renglones) debe ser breve; una respuesta para desarrollar debe desarrollarse. En ambas importa conocer lo que significa cada concepto.
    \item Lo que se escribe mal no se interpreta con benevolencia: quien lee no supone ``lo que se quiso decir''. Es necesario acostumbrarse tanto a la oralidad como a la escritura.
\end{itemize}

\section{Principios del derecho comunitario}

En la corrección se mencionaron los siguientes principios, cuyo significado debe conocerse:

\begin{itemize}
    \item \textbf{Autonomía.}
    \item \textbf{Especificidad.}
    \item \textbf{Supremacía:} supranacionalidad y supremacía son lo mismo, según lo expuesto.
    \item \textbf{Aplicabilidad inmediata,} con una excepción que no se distingue en las notas originales.
    % TODO: contenido incompleto o ambiguo en las notas originales (excepción a la aplicabilidad inmediata)
\end{itemize}

También se mencionaron, de forma fragmentaria, la armonización, la uniformidad y la armonía; el ciudadano europeo; y el conflicto de normas.
% TODO: contenido incompleto o ambiguo en las notas originales (conceptos del examen mencionados sin desarrollo)

\section{Cómo comparar dos sistemas}

Cuando se pide \textbf{comparar} no corresponde exponer todo un sistema (por ejemplo, el Mercosur) y a continuación todo el otro (la Unión Europea): eso no es comparación. Comparar es contrastar punto por punto (``esto es blanco, esto es negro''). Es válido trazar una línea en el medio y consignar cada característica de cada sistema; si queda tiempo, luego se desarrolla.

\begin{example}{Comparación del sistema de solución de controversias}
El sistema del Mercosur es \textbf{mixto}, porque tiene un mecanismo de heterocomposición y de autocomposición, e incluye un arbitraje ad hoc (ni siquiera permanente). El de la Unión Europea es una \textbf{instancia judicial} en la que se resuelve. Esta última es la instancia \textbf{obligatoria}; la del Mercosur es \textbf{alternativa}: se puede ir a ese mecanismo o a otro. Además, la cuestión prejudicial no es lo mismo que la opinión consultiva: en el Mercosur era opinión consultiva (si no se quiere, no se le hace caso); en la Unión Europea la prejudicial es obligatoria.
\end{example}

Un error frecuente es desarrollar todo el procedimiento del Mercosur cuando no es el tema pedido y limitarse a dos renglones sobre la Unión Europea (``el tribunal resuelve obligatoriamente''). Lo importante es que lo alternativo es el acceso. Las consignas pueden pedir desarrollar un sistema (por ejemplo, ``desarrolle el Tribunal Permanente de Revisión'' o ``desarrolle el sistema de la Unión Europea'') o compararlos, y cada tipo exige una respuesta distinta. Véase el cuadro comparativo del capítulo anterior.

% ------------------------------------------------------------
\chapter[La integración y el ciudadano]{La integración y el ciudadano}
% ------------------------------------------------------------

\section{Ciudadanía comunitaria y vida cotidiana}

Tras analizar la parte de derecho público de la integración se pasa al ámbito privado: algunos derechos se crean en el ámbito de la integración y surgen efectos en cada uno de los ciudadanos. La cuestión es cómo afectan estos derechos a los ciudadanos del Mercosur.

Quien posee una nacionalidad de un Estado europeo (por ejemplo, la italiana) posee lo que se denomina \textbf{ciudadanía comunitaria}. Quienes tienen también la ciudadanía europea reciben, cuando hay votaciones, la papeleta para votar.

En los orígenes los europeos percibieron la integración como una cuestión entre Estados, ``allá arriba''. Cuando comenzaron a objetar determinadas medidas, salieron a la calle, y actualmente vuelven a hacerlo, por otro tema, cuando consideran que el Estado se excede y se mete en cuestiones que tienen que ver con su vida particular.

La integración \textbf{afecta el día a día}, pero depende del \textbf{tipo de integración} de que se trate y de en qué ámbitos de la vida pueden inmiscuirse los órganos de integración (escala de áreas de integración).

\begin{important}
Los derechos creados en los ámbitos de integración afectan el día a día de cada ciudadano, pero solo en los ámbitos que el tipo de integración permite regular. Si la integración es comercial, un órgano de integración no puede regular, por ejemplo, la \textbf{restitución internacional de menores} ni la \textbf{adopción internacional}: sería una \textbf{autoridad incompetente}.
\end{important}

\section{Especificidad: lo que puede regular la integración}

\begin{definition}{Principio de especificidad}
Cada Estado parte transmite a los órganos de integración \textbf{facultades específicas}. Específico significa que el órgano solamente puede regular sobre esas materias.
\end{definition}

Las facultades se transmiten de manera específica y exclusiva: solo los órganos de integración pueden regular esas materias, y el Estado ya no puede hacerlo porque se las transfirió. Cuando los órganos de integración regulan ese tema, la norma afecta al ciudadano. Qué cuestiones se regulan depende del tipo de integración.

% ------------------------------------------------------------
\chapter[Derechos humanos en la UE]{Derechos humanos en la Unión Europea}
% ------------------------------------------------------------

\section{Origen del conflicto: las facultades de las comunidades}

Una de las primeras cuestiones que planteó conflicto en la Unión Europea fue el tema de los derechos humanos. En los orígenes, entre 1951 y 1957, se formaron las comunidades europeas:

\begin{itemize}
    \item \textbf{París, 1951:} Comunidad Europea del Carbón y del Acero (CECA).
    \item \textbf{Roma, 1957:} Comunidad Económica Europea (CEE) y Comunidad Europea de la Energía Atómica (Euratom).
\end{itemize}
% TODO: contenido incompleto o ambiguo en las notas originales (la docente enuncia París 51 con la CECA, Euratom y la CEE, y Roma del 57, y señala que una pregunta fue «capciosa» porque no se invertían los tratados)

Las facultades otorgadas a los órganos consistían en regular la comercialización, la producción, la investigación en carbón, acero y energía atómica, y las mercancías. La pregunta es si los órganos de integración tenían facultades para regular aspectos relativos a los derechos humanos; la respuesta inicial es que no.

\section{Primera etapa: el tribunal se declara incompetente}

En los primeros pasos, el tribunal comunitario (denominado en las notas ``Tribunal de la Unión Europea'') fue consultado en casos en que una normativa interna no estaba de acuerdo con alguna disposición comunitaria. Se le pidió la interpretación y el tribunal respondió que \textbf{no tenía competencia} para decidir sobre cuestiones de derechos humanos: el derecho comunitario es otra cosa y no tiene que ver con los derechos humanos. Posteriormente el tribunal se retractó.

El razonamiento del tribunal consideraba que las cuestiones de derechos humanos estaban \textbf{fuera del ámbito de las comunidades europeas}, de modo que no le correspondía decidir si una norma comunitaria y una norma de derechos humanos se contraponían. Era como sostener que el derecho comunitario es autónomo y que ambos planos no pueden cruzarse: ``van por carriles distintos''. Ese razonamiento estuvo equivocado.

Se advierte que la regulación de los contratos internacionales afecta a las personas, aunque no afecte directamente el derecho a transitar ni el derecho a establecer un negocio, y que al pensar en ``derechos humanos'' se piensa de inmediato en el derecho a la vida.

Existe también una convención europea, así como en América existe la Convención Americana: Europa tiene una \textbf{Convención Europea}.

\section{Generaciones de derechos}

\begin{definition}{Generaciones de derechos}
Los derechos humanos se agrupan en derechos de \textbf{primera}, \textbf{segunda} y \textbf{tercera generación}. Son categorías de derechos, y los más elementales son condición de los demás: si no existen, no puede ejercerse el derecho al voto (se necesita tener vida).
\end{definition}

Cuanto más alta es la categoría, se asume que las categorías más bajas están, en principio, aseguradas. La preocupación por el derecho al medio ambiente sano surge porque la alimentación y la vivienda ya están satisfechas: nadie se preocupa porque un colectivo tire humo si no puede llegar a fin de mes y comer. Del mismo modo, en un hogar no se paga el cable ni las plataformas si no se tiene dinero para la comida.

\begin{important}
Las generaciones de derechos se apoyan unas en otras: primero se asegura lo elemental (vida, libertad, alimentación, vivienda) y recién después se atienden derechos como el medio ambiente sano.
\end{important}

Cuando el tribunal aceptó el planteo, probablemente tuvo en consideración solamente los derechos humanos de \textbf{primera generación}: el derecho a la libertad y el derecho a la vida. Por eso afirmó que el derecho comunitario y los derechos humanos van por carriles distintos y que no tenía potestad para decidir si un derecho humano contradice o no el derecho comunitario.

\section{Debate: ¿se excede la especificidad?}

Se planteó la objeción de que el tribunal podría estar excediendo la especificidad: si las facultades se otorgaron para ciertas materias, no tendría por qué regularse sobre derechos humanos.

La respuesta es que \textbf{no se excede la especificidad} si se regula el derecho individual correspondiente a la materia delegada. Si se legisla sobre la capacidad de entrar y salir del área de integración (libre tránsito en un mercado común), sobre la libertad de trabajo o sobre el derecho a contratar, se regulan derechos individuales vinculados con la libre circulación de uno de los factores de producción: el trabajo.

\begin{important}
Regular el derecho individual correspondiente a la materia delegada no excede la especificidad (por ejemplo, libre tránsito, libertad de trabajo, derecho a contratar en un mercado común). Pero la regulación debe ser acorde ``pura y exclusivamente'' a eso: si las cárceles tienen que ser sanas y limpias, el Mercosur no se ocuparía, porque no se le delegaron esas facultades.
\end{important}

\section{Segunda y tercera etapa}

\subsection{Segunda etapa}

El tribunal consideró después que, en cierta medida, sí le corresponden las cuestiones de derechos humanos. Las constituciones de los Estados que integran la entonces Comunidad Económica Europea elaboraron como base los derechos fundamentales, en base a los derechos humanos; el respeto a los derechos humanos va en la naturaleza del \textbf{Estado de derecho}. Por lo tanto, cuando se delega una facultad específica, se la transmite con el presupuesto de que deben respetarse los derechos humanos. Aunque no se estableció específicamente, va de suyo.

Se encuadró la función de regular sobre esta materia en un derecho que se construye como presupuesto lógico a partir de los derechos humanos y del Estado de derecho: en la esencia de la facultad delegada está el respeto a ambos. Se buscó la manera de encuadrarlo aunque no estaba establecido específicamente en el tratado.

\subsection{Tercera etapa}

En una tercera etapa se incorpora la \textbf{Convención Europea} directamente al derecho comunitario. Cuando se habla de derechos humanos, se hace referencia a ellos tal como fueron reflejados en la Convención Europea.

\subsection{Evaluación del razonamiento inicial}

A juicio de la docente, el razonamiento inicial también incurrió en un error: la \textbf{libre circulación} y la \textbf{igualdad de trato} están desde los orígenes de las comunidades, de modo que ya había dos derechos individuales fundamentales de toda persona humana o jurídica, dos pilares en los que se construía el derecho europeo. El razonamiento contrario probablemente pensaba en el derecho a la vida y no en otros derechos, como el derecho al estudio. Para la docente era contradictorio, pues si existían disposiciones que regulaban la libertad de tránsito, no podía decirse que no correspondía porque iban por carriles distintos. Sin embargo, no existía una carta o declaración de derechos humanos suficientemente clara en el tratado.

\subsection{Ampliación de los derechos: desarrollo sustentable y medio ambiente}

Se fueron agregando derechos de primera, segunda y tercera generación. La regulación abarca al trabajador en actividad y también cuando cesa, la vejez digna y el medio ambiente sano, que sirve como factor de producción económica; por eso debe hacerse un \textbf{desarrollo sustentable} que garantice el medio ambiente. Ello significa no contaminar y buscar un desarrollo industrial con energías renovables, no contaminantes. Así, con la protección del medio ambiente, nace en cabeza de la persona el derecho al medio ambiente sano.

Se van agregando razonamientos que amplían los derechos sobre los cuales puede regular la Unión Europea. Es un camino que se recorrió desde la década de 1950 hasta la actualidad; no se produjo de la noche a la mañana.

\subsection{Cuadro de síntesis}

\begin{table}[H]
\centering
\small
\begin{tabularx}{\textwidth}{
    >{\raggedright\arraybackslash}p{0.14\textwidth}
    >{\raggedright\arraybackslash}X
    >{\raggedright\arraybackslash}X
}
\toprule
\textbf{Etapa} &
\textbf{Posición del tribunal} &
\textbf{Razonamiento} \\
\midrule

\textbf{Primera} &
Se declara incompetente: los derechos humanos quedan fuera del ámbito de las comunidades europeas. &
Derecho comunitario y derechos humanos ``van por carriles distintos''. Según la docente, quizá se pensó solo en los derechos de primera generación. \\

\textbf{Segunda} &
Admite que, en cierta medida, le corresponden esas cuestiones. &
Las constituciones de los Estados miembros se basan en derechos fundamentales; el respeto a los derechos humanos y al Estado de derecho es presupuesto de toda facultad delegada. \\

\textbf{Tercera} &
Incorpora la Convención Europea al derecho comunitario. &
Los derechos humanos se entienden tal como se reflejan en la Convención Europea; luego se amplían (desarrollo sustentable, medio ambiente sano). \\

\bottomrule
\end{tabularx}
\end{table}

% ------------------------------------------------------------
\chapter[Niveles de compromiso de los Estados]{Niveles de compromiso de los Estados}
% ------------------------------------------------------------

Cuando un Estado asume compromisos, sea de forma individual o a través de una organización, puede adquirir obligaciones de distinto nivel y con distinto compromiso. Los niveles tratados son los de \textbf{reconocimiento} (o respeto), \textbf{protección} y \textbf{promoción}.

\section{Reconocimiento y declaraciones}

\begin{definition}{Reconocimiento}
El Estado reconoce que la persona tiene un derecho (por ejemplo, a la vida, a una vivienda digna, a trabajar). Es un derecho de la persona y el Estado lo reconoce. La única obligación que asume es reconocerlo: nada más.
\end{definition}

Los instrumentos internacionales contienen \textbf{declaraciones} y \textbf{convenciones}. La declaración enuncia que ``todo individuo tiene derecho a\dots'': reconoce que la persona tiene tales derechos, pero no asume ningún compromiso. Las declaraciones se ubican en este nivel; las convenciones, en cambio, pueden asumir compromisos mayores.

\section{Declaraciones históricas}

\begin{itemize}
    \item \textbf{Declaración de la Revolución Francesa (1789):} la primera declaración de derechos humanos.
    \item \textbf{Declaración Universal de los Derechos Humanos (ONU):} diciembre de 1948. Es universal.
    \item \textbf{Declaración Americana de los Derechos y Obligaciones del Hombre (OEA):} mayo de 1948, anterior a la Declaración Universal; fue redactada el 30 y aprobada el 1.º de mayo. Es regional y es una declaración de derechos y obligaciones: solo la americana y la africana tienen un capítulo de obligaciones.
\end{itemize}
% TODO: contenido incompleto o ambiguo en las notas originales (fechas de redacción y aprobación de la Declaración Americana; en la exposición se corrigió de agosto a mayo; el ámbito institucional de cada declaración se menciona de forma confusa)

El primer documento a nivel general se dio, por lo tanto, en el ámbito americano. La Argentina ratificó la Declaración Americana, pero es una declaración: no asume compromisos. Se supone que quien reconoce un derecho a la vida lo va a respetar, pero no se obliga a realizar ningún acto positivo de cuidado o de respeto. Esto es derecho internacional (universal) y derecho americano (regional).

\section{Protección de derechos}

Cuando el Estado asume una obligación de \textbf{protección} se compromete a realizar un \textbf{acto concreto} para cuidar que la persona pueda ejercer el derecho. No basta con que se tenga derecho a la vida: si el Estado protege la vida, castiga a quien atente contra ella, dicta códigos que penan el homicidio y encarcela a los delincuentes; realiza actos para garantizar el ejercicio del derecho a la vida y de todos los demás derechos.

\begin{example}{Protección del derecho a la educación y a la libre circulación}
Si el Estado protege el derecho a la educación, no va a permitir las tomas de escuelas o universidades, pues obstaculizan el libre ejercicio de ese derecho; tendría que desalojar a quienes impiden estudiar. Si alguien corta una calle, impide el derecho a la libre circulación, y quien se comprometió a protegerlo debe despejar; se debe dejar un carril y las manifestaciones deben hacerse sobre la vereda, porque no pueden violar el derecho. El Estado protege el derecho mediante una reglamentación que garantice que se podrá circular y estudiar.
\end{example}

\section{Promoción de derechos}

El nivel de \textbf{promoción} consiste en \textbf{incentivar}: el Estado quiere que la persona ejerza el derecho, y no le da lo mismo que estudie o no. Quien protege garantiza que se pueda ir a la universidad; quien promueve crea universidades públicas y asigna becas, porque quiere que la persona ejerza ese derecho. Se compromete a realizar actos para fomentar el ejercicio del derecho.

Una cláusula en la que ``los Estados se comprometen a promover y desarrollar el derecho'' corresponde al nivel de promoción y no al de protección.

Por eso deben leerse con cuidado los tratados: muchas veces la fórmula es ``nos propondremos hacer tal cosa'', y debe indagarse qué se hará concretamente. Las declaraciones se ubican en este nivel.

\subsection{Cuadro de síntesis de los tres niveles}

\begin{table}[H]
\centering
\small
\begin{tabularx}{\textwidth}{
    >{\raggedright\arraybackslash}p{0.17\textwidth}
    >{\raggedright\arraybackslash}X
    >{\raggedright\arraybackslash}X
}
\toprule
\textbf{Nivel} &
\textbf{Qué asume el Estado} &
\textbf{Ejemplos dados} \\
\midrule

\textbf{Reconocimiento} &
Solo reconoce que la persona tiene el derecho. No asume ningún acto concreto. &
``Yo reconozco que vos tenés derecho a la vida / a una vivienda digna / a trabajar''. Las declaraciones. \\

\textbf{Protección} &
Realiza actos concretos para garantizar que la persona pueda ejercer el derecho. &
Códigos que penan el homicidio; no permitir tomas de escuelas o universidades; despejar calles cortadas; reglamentar para garantizar la circulación. \\

\textbf{Promoción} &
Incentiva y fomenta el ejercicio del derecho. &
Universidades públicas; asignación de becas. \\

\bottomrule
\end{tabularx}
\end{table}

\section{Convención Americana}

En una segunda etapa americana, hacia 1969, surge la \textbf{Convención Americana} o \textbf{Pacto de San José de Costa Rica}. Se crea un mecanismo de contralor de la OEA hacia los Estados, que pueden ser sancionados. Ya no se trata del mismo nivel: el Estado reconoció que va a poner en funcionamiento el derecho y, si no lo hace, es sancionado. Se menciona la Comisión de Derechos Humanos de la Convención Americana.

Como temas pendientes de desarrollo se anunciaron la Convención Americana y, fundamentalmente, la libertad de tránsito.

% ============================================================
% CORNELL
% ============================================================
\chapter[Cuadros Cornell]{Cuadros Cornell}

% ------------------------------------------------------------
\section{Cornell: aclaraciones conceptuales}

\begin{longtable}{
    >{\raggedright\arraybackslash}p{0.26\textwidth}
    >{\raggedright\arraybackslash}p{0.66\textwidth}
}
\toprule
\textbf{Preguntas / palabras clave} & \textbf{Notas / respuestas} \\
\midrule
\endhead

\textbf{¿Quién elabora, aprueba y controla el presupuesto en la Unión Europea?} &
El Ejecutivo (la Comisión, a través de los comisarios) elabora y ejecuta el plan; el Parlamento aprueba el presupuesto, lo controla y sanciona a los comisarios. Las preguntas formuladas con una acción (``¿quién legisla?'', ``¿quién juzga?'') llevan a identificar el órgano. \\

\textbf{¿Cómo se representa a la población y a los Estados?} &
El Parlamento representa a la población: hay un voto por representante y la composición varía según la proporcionalidad. En el Consejo hay un representante y un voto por cada Estado. \\

\textbf{¿Qué diferencia hay entre acceso alternativo y acceso obligatorio?} &
En el Mercosur (Tribunal Permanente de Revisión) el acceso es alternativo: puede elegirse otro foro, como la Organización Mundial de Comercio; la sentencia sí es obligatoria. En la Unión Europea el tribunal es el único que interpreta y aplica el derecho comunitario: el acceso es obligatorio y hay que acatar lo que diga. \\

\textbf{¿Qué distingue a la opinión consultiva de la cuestión prejudicial?} &
En el Mercosur la opinión consultiva puede no acatarse; en la Unión Europea la cuestión prejudicial es obligatoria. \\

\textbf{¿Qué es el derecho originario en el Mercosur?} &
El Tratado de Asunción es derecho originario porque crea el Mercosur. El Protocolo de Olivos también podría considerarse originario, pero no creó el Mercosur. No debe confundirse con el ``Pacto de Olivos'', que es ``la reunión'' y no tiene relación con el tema. \\

\midrule
\multicolumn{2}{p{\dimexpr0.92\textwidth+2\tabcolsep\relax}}{
\textbf{Resumen:}
El presupuesto europeo se elabora y ejecuta por el Ejecutivo (la Comisión) y se aprueba y controla por el Parlamento, que representa a la población, mientras que en el Consejo hay un voto por Estado. El acceso al tribunal es alternativo en el Mercosur (la sentencia sí es obligatoria) y obligatorio en la Unión Europea. El Tratado de Asunción es derecho originario porque crea el Mercosur.
} \\
\bottomrule
\end{longtable}

% ------------------------------------------------------------
\section{Cornell: criterios de estudio y comparación}

\begin{longtable}{
    >{\raggedright\arraybackslash}p{0.26\textwidth}
    >{\raggedright\arraybackslash}p{0.66\textwidth}
}
\toprule
\textbf{Preguntas / palabras clave} & \textbf{Notas / respuestas} \\
\midrule
\endhead

\textbf{¿Qué debía consignarse en el cuadro por tratado fundacional?} &
El área de integración y sus características; si es intergubernamental o supranacional; si es unión aduanera o zona de libre comercio; y si las decisiones se resuelven por consenso o por órganos supranacionales. \\

\textbf{¿Por qué hay que leer bien la consigna?} &
Porque puede haber más de una opción correcta, hay que cuidar los términos (alternativo se refiere al acceso) y la extensión de la respuesta debe ajustarse a lo pedido (breve o para desarrollar). \\

\textbf{¿Qué principios del derecho comunitario se mencionaron?} &
Autonomía, especificidad, supremacía (supranacionalidad y supremacía son lo mismo, según lo expuesto) y aplicabilidad inmediata. Importa saber qué significa cada uno. \\

\textbf{¿Cómo se comparan correctamente dos sistemas?} &
No se describe uno y después el otro: se contrastan punto por punto, por ejemplo en una tabla. Mercosur: sistema mixto (autocomposición y heterocomposición, arbitraje ad hoc), acceso alternativo, opinión consultiva. Unión Europea: sistema judicial, acceso obligatorio, cuestión prejudicial obligatoria. \\

\midrule
\multicolumn{2}{p{\dimexpr0.92\textwidth+2\tabcolsep\relax}}{
\textbf{Resumen:}
Para responder bien hay que identificar los criterios del estudio de cada tratado fundacional, leer con atención la consigna y cuidar los términos. Comparar es contrastar punto por punto; el sistema del Mercosur es mixto y de acceso alternativo, y el de la Unión Europea es judicial y de acceso obligatorio.
} \\
\bottomrule
\end{longtable}

% ------------------------------------------------------------
\section{Cornell: la integración y el ciudadano}

\begin{longtable}{
    >{\raggedright\arraybackslash}p{0.26\textwidth}
    >{\raggedright\arraybackslash}p{0.66\textwidth}
}
\toprule
\textbf{Preguntas / palabras clave} & \textbf{Notas / respuestas} \\
\midrule
\endhead

\textbf{¿Cómo afecta la integración a cada ciudadano?} &
Los derechos creados en el ámbito de integración afectan el día a día, según el tipo de integración y los ámbitos que los órganos pueden regular. Ejemplo: quien tiene ciudadanía comunitaria (la italiana) recibe la papeleta para votar. \\

\textbf{¿Qué es la especificidad?} &
Los Estados transmiten a los órganos de integración facultades específicas y exclusivas: solo pueden regular sobre esas materias y el Estado ya no puede hacerlo. \\

\textbf{¿Qué ocurre si un órgano regula fuera de lo delegado?} &
Es una autoridad incompetente. Ejemplo: en una integración comercial, un órgano no puede regular la restitución internacional de menores ni la adopción internacional. \\

\midrule
\multicolumn{2}{p{\dimexpr0.92\textwidth+2\tabcolsep\relax}}{
\textbf{Resumen:}
La integración afecta a los ciudadanos solo en las materias que los Estados delegaron de manera específica y exclusiva; fuera de ellas, el órgano de integración es incompetente.
} \\
\bottomrule
\end{longtable}

% ------------------------------------------------------------
\section{Cornell: derechos humanos en la Unión Europea}

\begin{longtable}{
    >{\raggedright\arraybackslash}p{0.26\textwidth}
    >{\raggedright\arraybackslash}p{0.66\textwidth}
}
\toprule
\textbf{Preguntas / palabras clave} & \textbf{Notas / respuestas} \\
\midrule
\endhead

\textbf{¿Podían las comunidades regular derechos humanos en su origen?} &
No: sus facultades se limitaban a la comercialización, la producción, la investigación en carbón, acero y energía atómica, y las mercancías (CECA, 1951; CEE y Euratom, 1957). \\

\textbf{¿Por qué el tribunal europeo rechazó primero los derechos humanos?} &
Se declaró incompetente porque entendió que estaban fuera del ámbito de las comunidades: derecho comunitario y derechos humanos ``van por carriles distintos''. \\

\textbf{¿Cómo cambió la postura del tribunal?} &
En la segunda etapa reconoció que el respeto a los derechos humanos y al Estado de derecho es presupuesto de toda facultad delegada. En la tercera incorporó la Convención Europea y amplió los derechos (desarrollo sustentable, medio ambiente sano). \\

\textbf{¿Qué son las generaciones de derechos?} &
Categorías de derechos (primera, segunda y tercera generación) que se apoyan unas en otras: primero lo elemental (vida, libertad, alimentación, vivienda) y luego los derechos más amplios (medio ambiente sano). Ejemplo: no se paga el cable si no hay para comer. \\

\textbf{¿Se excede la especificidad al regular derechos humanos?} &
No, si se regula el derecho individual correspondiente a la materia delegada (libre tránsito, libertad de trabajo, derecho a contratar). Sí se excedería con materias no delegadas, como las condiciones de las cárceles. \\

\midrule
\multicolumn{2}{p{\dimexpr0.92\textwidth+2\tabcolsep\relax}}{
\textbf{Resumen:}
El tribunal europeo pasó de declararse incompetente en derechos humanos a incorporarlos como presupuesto del Estado de derecho y mediante la Convención Europea, con una ampliación progresiva hacia el desarrollo sustentable. Regular el derecho individual correspondiente a la materia delegada no excede la especificidad.
} \\
\bottomrule
\end{longtable}

% ------------------------------------------------------------
\section{Cornell: niveles de compromiso de los Estados}

\begin{longtable}{
    >{\raggedright\arraybackslash}p{0.26\textwidth}
    >{\raggedright\arraybackslash}p{0.66\textwidth}
}
\toprule
\textbf{Preguntas / palabras clave} & \textbf{Notas / respuestas} \\
\midrule
\endhead

\textbf{¿Qué obliga a un Estado una declaración?} &
Solo reconoce que la persona tiene el derecho; no asume ningún compromiso ni acto positivo. \\

\textbf{¿Cuáles son los tres niveles de compromiso de un Estado?} &
Reconocimiento: solo declara el derecho. Protección: actos concretos para garantizar su ejercicio (penar el homicidio, despejar calles, impedir tomas). Promoción: incentivar y fomentar (universidades públicas, becas). \\

\textbf{¿En qué se diferencia proteger de promover?} &
Proteger es garantizar que la persona pueda ejercer el derecho; promover es incentivar activamente que lo ejerza. \\

\textbf{¿Qué declaraciones y convenciones se mencionaron?} &
Revolución Francesa (1789); Declaración Americana de los Derechos y Obligaciones del Hombre (mayo de 1948, regional, anterior a la universal); Declaración Universal (diciembre de 1948); Convención Americana o Pacto de San José (1969), con mecanismo de control de la OEA y sanción al Estado. \\

\midrule
\multicolumn{2}{p{\dimexpr0.92\textwidth+2\tabcolsep\relax}}{
\textbf{Resumen:}
Los Estados pueden asumir obligaciones de reconocimiento, protección o promoción. Las declaraciones se limitan a reconocer; la protección exige actos concretos; la promoción, incentivar el ejercicio del derecho. La Convención Americana de 1969 introduce un mecanismo de control y sanción.
} \\
\bottomrule
\end{longtable}

\end{document}
